\section{Introduction}

Designing a protein that displays a target epitope in its native conformation is a route
to immunogens and binders that recapitulate a known antibody contact. We test the
RFdiffusion3 $\rightarrow$ ProteinMPNN $\rightarrow$ AlphaFold3 pipeline against two
questions.

\paragraph{Backbone geometry vs.\ sequence distribution (DP3 set).}
An RFdiffusion3-only pipeline passes AlphaFold3 validation far less often than Lawson's
RFdiffusion1 + ProteinMPNN baseline. The competing explanations are that RFdiffusion3
produces poorer \emph{backbone geometry}, or that its \emph{sequences} are
out-of-distribution for AlphaFold3. Re-sequencing the RFdiffusion3 backbones with
ProteinMPNN (trained on the inverse-folding task, sharing AlphaFold3's natural-PDB
distribution) isolates the two: if RFD3+MPNN approaches the baseline pass rate, the
sequence-distribution hypothesis is supported. The four-filter definitions and run
layout are recorded in \texttt{docs/README\_v2\_original.md}.

\paragraph{Scaffolding without an antibody (12-mer set).}
Many targets have no known antibody, so there is no antibody-clash filter. We instead
scaffold tiled 12-residue epitopes and judge antibody \emph{accessibility} with a
native-antigen-aware cylinder surrogate (Section~\ref{sec:methods}): it asks whether
non-epitope scaffold residues intrude into the volume an approaching antibody would
occupy, while excusing residues that merely sit where the native antigen already sits.

A single config-driven composite scorer (Section~\ref{sec:methods}) ranks designs in
both settings, differing only by which accessibility term it uses.
